\section{Completing the seven cohomological fields on one linked frame}
\label{sec:linked-frame-completion}

\paragraph{The endpoint.}
This chapter proves all seven fields left open by MFC14:
\texttt{phi}, \texttt{gen}, \texttt{cl}, \texttt{ec},
\texttt{b1p}, \texttt{b1}, and \texttt{h1}.
The proofs apply to the actual linked frame already constructed there.
They neither replace its quotient nor enlarge its class-rank constant.
Consequently the original \texttt{frame\_ex} theorem, with its full
\texttt{Gen} and \(\texttt{GC.Hypotheses}_{2}\) antecedents, is obtained.
For the original coefficient DVR, CAF7 then installs the character
\texttt{SeedClause} and gives \texttt{HasFrame}.
The separate \texttt{gen\_gd} and \texttt{gen\_host} entry theorems and
the two automorphic seed representations are not asserted here.

\paragraph{Retained mathematical inputs.}
Write \(P:\texttt{GC.Hypotheses}_{2}\) for the original published
cohomological package.  We use its local finiteness, local Tate
duality, local Euler characteristic, global finiteness, global Euler
characteristic and normalized Poitou--Tate fields at their stated
finite-module instances.  MFC's two strict classes retain the
Greenberg--Wiles input used there.  Class-group finiteness, the
Dirichlet unit theorem and finite class-field reciprocity are the
named classical number-theoretic inputs below; the last is the
same class-character dictionary already used in MFC5 and TW1.
We give the needed Kummer, coefficient, quotient, inertia and cone
arguments explicitly.  In particular the specialized source
\texttt{CharacterH1GaloisQuotient} theorem is replaced, for its
trivial-character application here, by a CM-unit argument.

\subsection{The linked data and literal cohomology spaces}

Until the final assembly let \(F\) be a totally real number field,
\(\Omega\) its specified algebraic closure, and
\(G_F=\operatorname{Gal}(\Omega/F)\).  Let \(\ell\) be odd and
\(k\) a field of characteristic \(\ell\).  The given frame
\(\Phi\) has a profinite group \(\Gamma\), a sign character
\(\delta:\Gamma\to\{\pm1\}\), an element \(c\) with
\(\delta(c)=-1\), and \(\Gamma_0=\ker\delta\).
It has finite sets \(V,\mathrm{Inf}\), profinite groups \(G_v\),
maps
\[
 \operatorname{loc}_{0,v}:G_v\longrightarrow\Gamma_0,\qquad
 \operatorname{loc}_v:G_v\longrightarrow\Gamma,
 \qquad I_v\subset G_v,\qquad d_v\in\mathbb N,
\]
the cyclotomic character \(\epsilon:\Gamma\to\mathbb Z_\ell^\times\),
and a residual character \(\Psi:\Gamma\to k^\times\).
Put
\[
 \operatorname{loc}_{2,(v,0)}=\operatorname{loc}_{0,v},\qquad
 \operatorname{loc}_{2,(v,1)}(s)
       =c\,\operatorname{loc}_{0,v}(s)c^{-1},\qquad
 N_v=I_v\cap\ker(\epsilon\operatorname{loc}_v).
\]
The element in the second formula is regarded as an element of
\(\Gamma_0\).  For any characteristic-\(\ell\) field \(E\), denote
the reduction of \(\epsilon\) in \(E^\times\) by \(\omega_E\),
and the sign character in \(E^\times\) by \(\delta_E\).

The full \texttt{FLink} gives a finite set \(S\) consisting
\emph{exactly} of the places of \(F\) above \(\ell\), a normal open
subgroup \(K\subset G_F\) of index prime to \(\ell\), and a
continuous surjection
\[
 \pi:G_F\twoheadrightarrow\Gamma,\qquad
 \ker\pi=
 \overline{\langle I_A\cap K:
                  A\mid v,\ v\notin S\rangle^{\mathrm{normal}}}.
                                                        \tag{LFC.1}
\]
There are bijectively indexed places \(v\mapsto\mathfrak p_v\in S\),
chosen prolongations \(A_v\), and continuous isomorphisms
\(\iota_v:G_v\simeq D_{A_v}\).  Under these isomorphisms,
\(\operatorname{loc}_v=\pi\iota_v\), \(I_v\) is the actual inertia
group, and \(d_v=[F_{\mathfrak p_v}:\mathbb Q_\ell]\).
The set \(\mathrm{Inf}\) indexes all real places, and \(c_w\)
is the image of a complex conjugation at that place.
Furthermore
\[
 \epsilon\pi=\epsilon_F,\qquad
 K\subset\ker(\delta\pi)\cap\ker(\Psi\pi),\qquad
 \ker\delta\ \hbox{is open}.                              \tag{LFC.2}
\]
No assertion of \texttt{Ok} is included in this input.
When only \texttt{GLink} is needed, the hypotheses are (LFC.1),
\(\ell\nmid[G_F:K]\), and \(S\) containing the \(\ell\)-places.

All cochains below are continuous, with discrete coefficient
modules.  For a representation \(\rho\) on \(W\), write
\[
 (d x)(g)=\rho(g)x-x,\qquad
 (d z)(g,h)=z(g)+\rho(g)z(h)-z(gh).
\]
For stable local subspaces \(b_v\subset W\), the source's
\(H^1_C(\operatorname{loc},W,b)\) consists of pairs
\((z,(x_v))\), where \(z\) is a cocycle and
\[
 z(\operatorname{loc}_v s)
       -\bigl(\rho(\operatorname{loc}_v s)x_v-x_v\bigr)\in b_v,
                                                        \tag{LFC.3}
\]
modulo the pairs \((dX,(X+y_v))\), \(y_v\in b_v\).
This is the literal \texttt{Z1C}/\texttt{continuousCoboundaries1}
quotient, including its local primitives.  In particular local
condition \(0\) retains those primitives, whereas local condition
\(W\) makes them irrelevant.

\subsection{Finite characters and the uniform class exponent}

% BEGIN LFC1
\paragraph{LFC1. Kummer classes with bounded ramification are finite.}
\label{claim:LFC1}
Let \(E\) be a number field containing \(\mu_\ell\), and let
\(T\) be a finite set of its finite places.  Put
\[
 V_T(E)=\{a\in E^\times:\ell\mid v(a)\text{ for every }v\notin T\}.
\]
There is an exact sequence of finite-dimensional
\(\mathbb F_\ell\)-spaces
\[
 0\longrightarrow
   \mathcal O_{E,T}^{\times}/(\mathcal O_{E,T}^{\times})^\ell
 \longrightarrow V_T(E)/E^{\times\ell}
 \longrightarrow \operatorname{Cl}(\mathcal O_{E,T})[\ell]
 \longrightarrow0.                                     \tag{LFC.4}
\]
Consequently there are only finitely many continuous additive
characters \(G_E\to\mathbb F_\ell\) unramified outside \(T\).

\begin{proof}
In the ring with the primes \(T\) inverted, write
\((a)=\mathfrak a^\ell\).  Unique ideal factorization makes
\(\mathfrak a\) unique, and sends the class of \(a\) to
\([\mathfrak a]\).  Multiplying \(a\) by an \(\ell\)-th power
changes \(\mathfrak a\) by a principal ideal.  Conversely an ideal
class killed by \(\ell\) has a representative with
\(\mathfrak a^\ell=(a)\), giving surjectivity.  The class is zero
precisely when \(a=u b^\ell\), with \(u\) a \(T\)-unit.
If a \(T\)-unit is an \(\ell\)-th power in \(E^\times\),
its root is again a \(T\)-unit by valuations.  This proves
exactness.  The unit theorem and class-group finiteness prove
the stated finiteness.

Fix a primitive root \(\zeta\in E\).  A nonzero character
\(\chi:G_E\to\mathbb F_\ell\) cuts out a cyclic extension of
degree \(\ell\).  If \(\tau\) generates its Galois group with
\(\chi(\tau)=1\), the operator
\(\sum_{j=0}^{\ell-1}\zeta^{-j}\tau^j\) is nonzero by linear
independence of distinct field embeddings.  Choose \(x\) on
which it is nonzero and call its value \(\beta\).
Then \(\tau\beta=\zeta\beta\), \(\beta^\ell=a\in E^\times\),
and \(g\beta=\zeta^{\chi(g)}\beta\).
The zero character corresponds to \(a=1\).  If two such
\(a\)'s differ by an \(\ell\)-th power, the corresponding
\(\beta\)'s differ by an element of \(E^\times\) times a power
of \(\zeta\), so the two characters are equal.

At a place outside \(T\), the extension cut out by \(\chi\)
is unramified.  Its normalized valuation therefore takes integral
values on \(\beta\), whence \(v(a)=\ell v(\beta)\) is divisible
by \(\ell\).  The characters inject into the finite middle
space of (LFC.4).  This proves the final assertion without
assuming the converse local ramification criterion.
\end{proof}
% END LFC1

% BEGIN LFC2
\paragraph{LFC2. The full \texttt{PhiCond} on every open subgroup.}
\label{claim:LFC2}
For any \texttt{GLink} over any number field \(F\), without
a total-reality assumption, every open subgroup \(H\subset\Gamma\)
has only finitely many homomorphisms \(H\to(\mathbb F_\ell,+)\)
with open kernel.  Thus \(\operatorname{PhiCond}(\Gamma,\ell)\)
holds, including its quantifier over all open \(H\).

\begin{proof}
Let
\[
 J=\pi^{-1}(H)\cap K\cap\ker(\bar\epsilon_F),\qquad
 E=\Omega^J.
\]
This is an open subgroup, so \(E/F\) is finite, and \(E\)
contains \(\mu_\ell\).  The map \(G_E=J\to H\) induced by
\(\pi\) has image of finite index: representatives of the
finitely many cosets of \(J\) in \(G_F\) also cover its image
cosets in \(\Gamma\), and hence in \(H\).
Choose a nonzero rational integer \(n\) divisible by the residue
characteristic of every place of \(S\).
At a place of \(E\) not dividing \(n\), the underlying place of
\(F\) is outside \(S\).  Its inertia in \(G_E\) is contained
in \(I_A\cap K\), and is killed by (LFC.1).

Consequently restriction sends each character of \(H\) to one
of the finitely many characters supplied by LFC1, with \(T\)
the places of \(E\) dividing \(n\).  Fix finitely many coset
representatives for the image of \(G_E\) in \(H\).  A character
is determined by its restriction and its value on these
representatives, so there are only finitely many possibilities.
Normality of this image is not needed for this argument.
For a homomorphism to a finite discrete group, an open kernel
is equivalent to continuity, since every fiber is a coset of
the kernel.  This is exactly the source's definition of
\texttt{PhiCond}, not a statement only about \(\Gamma\) itself.
\end{proof}
% END LFC2

% BEGIN LFC3
\paragraph{LFC3. One nonzero exponent for all unramified characters.}
\label{claim:LFC3}
For a full \texttt{FLink}, there is \(h>0\), chosen before
the finite abelian \(\ell\)-group \(B\) and the character
\(\chi:\Gamma_0\to B\), such that every open-kernel \(\chi\)
satisfying
\[
 \chi(\operatorname{loc}_{2,(v,b)}s)=1
 \quad(v\in V,\ b\in\{0,1\},\ s\in I_v)
                                                        \tag{LFC.5}
\]
satisfies \(\chi(g)^h=1\) for every \(g\in\Gamma_0\).
This is the full \texttt{cl} field.

\begin{proof}
Put \(U=\pi^{-1}(\Gamma_0)\), \(E=\Omega^U\), and choose
\[
                  h=2\,|\operatorname{Cl}(\mathcal O_E)|.
                                                        \tag{LFC.6}
\]
This positive integer depends only on the linked data.
The pullback of \(\chi\) to \(G_E=U\) is unramified at every
finite place.  Indeed, above \(v\in S\), conjugate the
prolongation to the chosen \(A_v\).
If the image \(a\) of the conjugating element has
\(\delta(a)=1\), the character sees a conjugate inside
\(\Gamma_0\) of \(\operatorname{loc}_{0,v}(I_v)\).
If \(\delta(a)=-1\), write \(a=(ac^{-1})c\);
it sees a conjugate inside \(\Gamma_0\) of the second branch
\(c\operatorname{loc}_{0,v}(I_v)c^{-1}\).
Both are killed by (LFC.5), since \(B\) is abelian.
This uses the actual inertia groups, including their
intersections with \(G_E\).

At a finite place outside \(S\), let \(t\) be such an inertia
element and put \(r=[G_F:K]\).  Normality of \(K\) gives
\(t^r\in K\), so (LFC.1) gives \(\pi(t)^r=1\).
The order of \(\chi(\pi(t))\) is an \(\ell\)-power and divides
the prime-to-\(\ell\) integer \(r\); hence it is one.
Finally the square of the pullback character kills every
real complex conjugation, since these have order at most two.
Finite class-field reciprocity therefore factors that square
through the ordinary finite ideal class group of \(E\).
It is annihilated by the class number, giving (LFC.6).
The map \(G_E\to\Gamma_0\) is onto, by surjectivity of \(\pi\),
so this proves the assertion for every \(g\).
No bound on \(B\), or choice depending on \(\chi\), entered
the choice of \(h\).
\end{proof}
% END LFC3

\subsection{Local generators with the required cyclotomic exception}

% BEGIN LFC4
\paragraph{LFC4. The elementary procyclic cohomology calculation.}
\label{claim:LFC4}
In this claim \(\ell\) may be any prime and the base number
field need not be totally real.
Let \(D\) be a local decomposition group, \(I\) its inertia,
and \(\varphi\in D\) a lift of arithmetic Frobenius, so that
\(D/I\simeq\widehat{\mathbb Z}\), with \(\varphi\) mapping to
one.  If \(J\triangleleft D\) is closed, contained in \(I\),
and \([I:J]\) is finite and prime to \(\ell\), then for every
finite-dimensional finite-field characteristic-\(\ell\)
continuous representation \(W\) on which \(J\) acts trivially,
\[
 H^2(D/J,W)=0,\qquad
 \dim H^1(D/J,W)=\dim H^0(D/J,W).                         \tag{LFC.7}
\]

\begin{proof}
For a procyclic group with generator \(t\), a continuous
one-cocycle is determined by \(x=z(t)\).
Every \(x\in W\) occurs: if the action has period \(a\),
the cocycle on the integers defined by its partial sums
factors through the cyclic group of order \(a\ell\), because
the sum over \(a\ell\) terms is zero.
Coboundaries correspond to \((t-1)W\).  Thus
\(H^1=W/(t-1)W\), whose dimension is \(\dim\ker(t-1)\).
A continuous two-cocycle factors through some finite cyclic
quotient.  For a cyclic quotient of order \(a\) acting through
\(t\), the periodic resolution gives
\(H^2=W^t/(1+t+\cdots+t^{a-1})W\): this follows by lifting
the cyclic generator in the extension defined by the cocycle,
whose \(a\)-th power is the displayed invariant.
Inflation to the quotient of order \(a\ell\) multiplies that
invariant by \(\ell\), and hence kills the class.
The same conclusion holds in continuous cohomology, since
the new primitive is a cochain on a finite quotient.

Now \(A=I/J\) is finite of order invertible in the coefficient
field.  Averaging over \(A\) projects \(W\) onto \(W^A\).
It also contracts the positive-degree homogeneous cochains
of \(A\): insert an initial variable \(a\), sum over \(a\in A\),
and divide by \(|A|\).  The cochain identity is \(ds+sd=1\)
in positive degrees.  Apply this contraction to the
normal-subgroup double complex for
\(1\to A\to D/J\to\widehat{\mathbb Z}\to1\);
all positive \(A\)-rows disappear and the remaining row is
the continuous cochain complex of
\(\widehat{\mathbb Z}\) with coefficients \(W^A\).
This proves (LFC.7) by the preceding calculation.
All maps are continuous: the \(A\)-sums are finite and every
cochain used here factors through a finite quotient.
\end{proof}
% END LFC4

% BEGIN LFC5
\paragraph{LFC5. Exactly \(d\) universal inertia detectors.}
\label{claim:LFC5}
Let \(v\mid\ell\), \(D=D_A\), \(I=I_A\), and
\(d=[F_v:\mathbb Q_\ell]>0\).  Assume the local fields of
\(P\), and assume \(\bar\epsilon_F|_D\ne1\).
For any normal open \(K\subset G_F\) of prime-to-\(\ell\)
index there are \(u_1,\ldots,u_d\in I\cap K\), with a
distinguished index \(a_0\), such that
\[
 \epsilon_F(u_a)=1\quad(a\ne a_0),                       \tag{LFC.8}
\]
and, for every finite abelian \(\ell\)-group \(B\), every
open-kernel homomorphism \(\chi:D\to B\) which kills all
\(u_a\) kills all of \(I\).

\begin{proof}
The residue-field Galois group is \(\widehat{\mathbb Z}\):
on its finite subfields Frobenius has precisely their degrees
as orders, and taking the inverse limit gives this statement.
Choose a Frobenius lift \(\varphi\in D\).  Then \(I\) and
\(\varphi\) topologically generate \(D\), and for every
integer \(n>0\) there is a continuous degree map
\(\deg_n:D\to\mathbb Z/n\) killing \(I\) and taking
\(\varphi\) to one.

Let \(X=\operatorname{Hom}_{\rm cont}(D,\mathbb F_\ell)\).
Local finiteness makes \(X\) finite-dimensional.
Local Tate duality gives \(H^2(D,\mathbb F_\ell)=0\):
its dual is the invariant subspace of the nontrivial
cyclotomic character.  The local Euler characteristic therefore
gives \(\dim X=d+1\).
Evaluation at \(\varphi\) is onto, witnessed by \(\deg_\ell\).
Thus \(Y=\{f\in X:f(\varphi)=0\}\) has dimension \(d\).
Put \(J=I\cap K\), \(J_0=J\cap\ker\epsilon_F\), and
\[
                   W_0=\{f\in Y:f|_{J_0}=0\}.
 \]
Evaluation on \(J\) separates \(Y\).  Indeed, if \(f|_J=0\)
and \(r=[G_F:K]\), then \(t^r\in J\) for \(t\in I\);
as \(r\) is invertible modulo \(\ell\), \(f|_I=0\).
Together with \(f(\varphi)=0\) and continuity this gives
\(f=0\).

If \(W_0=0\), successively choose points of \(J_0\) which
reduce the common evaluation kernel.  After at most \(d\)
choices that kernel is zero.  Pad the list by identities.
If \(W_0\ne0\), the same procedure chooses at most \(d-1\)
points of \(J_0\) with common kernel \(W_0\).
Let \(N\) be the common kernel in \(J\) of the finitely many
characters in \(W_0\).  It is open in the compact group \(J\)
and contains \(\ker(\epsilon_F|_J)\).
Hence \(J/N\) is a finite continuous quotient of the closed
subgroup \(\epsilon_F(J)\subset\mathbb Z_\ell^\times\).
Such a quotient is cyclic when \(\ell\) is odd.
For completeness, compactness and openness show that its
kernel contains the intersection with \(1+\ell^a\mathbb Z_\ell\)
for some \(a\): otherwise a nested sequence outside the
kernel would have a limit in the kernel with cyclotomic value
one.  The quotient is therefore a quotient of a subgroup of
\((\mathbb Z/\ell^a)^\times\), a cyclic group for odd \(\ell\).
Choose \(u_{a_0}\in J\) mapping to a generator of \(J/N\).
Vanishing at this one point forces a member of \(W_0\) to
vanish on \(J\), and hence to be zero.  Pad with identities
to obtain \(d\) points.  In both cases (LFC.8) holds, and
a character in \(X\) killed by the \(u_a\)'s and \(\varphi\)
is zero.

To obtain the assertion for every \(B\), put \(n=|B|\) and
consider the finite abelian \(\ell\)-group
\[
             C=\operatorname{im}(\chi,\deg_n)
                       \subset B\times\mathbb Z/n.
 \]
It is generated by the image of \(\varphi\).
Otherwise its quotient by that cyclic subgroup is a nonzero
finite abelian \(\ell\)-group and has a nonzero homomorphism
to \(\mathbb F_\ell\) (reduce one cyclic factor modulo \(\ell\)).
Composing with \(D\to C\) contradicts the preceding detection
statement: it kills all \(u_a\) and \(\varphi\).
For \(t\in I\), write its image in \(C\) as the \(b\)-th
power of the image of \(\varphi\).  Its degree is zero, so
\(n\mid b\); since \(n\) annihilates \(B\), \(\chi(t)=1\).
This also handles \(B=1\).  The same \(u_a\)'s work for all
groups \(B\) and all characters \(\chi\).
\end{proof}
% END LFC5

% BEGIN LFC6
\paragraph{LFC6. The exact local \texttt{gen} field.}
\label{claim:LFC6}
Assume \(P\), a full \texttt{FLink}, and
\(\omega_k\operatorname{loc}_v\ne1\) for every \(v\).
Then \(\Phi\) satisfies the full \texttt{gen} field:
for every \(v\) there are \(u:\operatorname{Fin}(d_v)\to G_v\)
and \(a_0\in\operatorname{Fin}(d_v)\), with
\[
 u_a\in I_v,\quad \Psi(\operatorname{loc}_v u_a)=1,\quad
 a\ne a_0\Longrightarrow u_a\in N_v,
                                                        \tag{LFC.9}
\]
which detect the restriction to inertia of every continuous
character to every finite abelian \(\ell\)-group.

\begin{proof}
The inclusion \(\mathbb F_\ell\hookrightarrow k\) is injective.
The compatibility \(\epsilon\pi=\epsilon_F\) thus identifies
the stated nontriviality with
\(\bar\epsilon_F|_{D_{A_v}}\ne1\), including the source's
modulo-\(\ell^1\) convention.
Apply LFC5 at \(A_v\) and pull its points back through the
continuous isomorphism \(\iota_v:G_v\simeq D_{A_v}\).
Every point lies in \(I_{A_v}\cap K\), so its image has
\(\Psi=1\) by (LFC.2).  All but the distinguished point
have cyclotomic value one, giving membership in \(N_v\).
The local degree is exactly \(d_v\), not an upper bound.
A character of \(G_v\) transfers through \(\iota_v^{-1}\)
to a character of \(D_{A_v}\) with open kernel, so LFC5
proves its universal inertia assertion verbatim.
\end{proof}
% END LFC6

\subsection{The actual strict-cohomology dictionary and rank bounds}

% BEGIN LFC7
\paragraph{LFC7. Literal class characters and extension of coefficients.}
\label{claim:LFC7}
Here \(F\) may be any number field and \(\ell\) any prime.
For an open subgroup \(J\subset G_F\), let
\[
 X_E(J)=\{f:J\to E:\ f\text{ locally constant and additive},
                   \ f(I_A\cap J)=0\text{ for every }A\}.
 \]
This is exactly \(\texttt{Kum.XDual}(1,J,E,E)\).
For every characteristic-\(\ell\) field \(E\),
\[
 \dim_E X_E(J)=\dim_{\mathbb F_\ell}X_{\mathbb F_\ell}(J)
                  =\operatorname{clRk}_\ell(J).           \tag{LFC.10}
\]
In particular the literal \texttt{clRk}, not a substitute
class number, bounds these coefficient spaces.

\begin{proof}
The extra decomposition condition in \texttt{XDual} is
imposed only when \(A(1)<1\), which never happens.
The other three conditions in its definition are exactly
those displayed above.  Identify \(J=G_{E_0}\) for the finite
fixed field \(E_0\).  Inertia under this identification is
the intersection with \(J\), by its definition as the
kernel of the action on the residue field.
Finite class-field reciprocity identifies the prime-field
space with the characters of the narrow ideal class group.
For odd \(\ell\) these are also the characters of the
ordinary class group, since the extra sign quotient is
a group of order a power of two.
Thus this prime-field space is finite-dimensional.

Any \(f\in X_E(J)\) has finite image, because \(J\) is
compact and \(f\) is locally constant.  Its image spans a
finite-dimensional \(\mathbb F_\ell\)-subspace of \(E\).
Take coordinates in a basis of that span.  Each coordinate
is an additive, locally constant function killing every
inertia intersection, and hence belongs to
\(X_{\mathbb F_\ell}(J)\).
It follows that a prime-field basis of that space spans
\(X_E(J)\) over \(E\).  It remains independent: finitely
many evaluations give a nonsingular minor over
\(\mathbb F_\ell\), which remains nonsingular over \(E\).
This proves (LFC.10), without assuming that \(E\) is finite.
\end{proof}
% END LFC7

% BEGIN LFC8
\paragraph{LFC8. The full strict-generator estimate.}
\label{claim:LFC8}
Here again \(F\) is any number field and \(\ell\) any prime.
Let \(U\subset G_F\), let \(J\triangleleft G_F\) be open with
\(J\subset U\) and \(\ell\nmid[U:J]\), and let
\(\pi_U:U\twoheadrightarrow\Delta\) be continuous.
Here \(\Delta\) is an abstract group equipped with any topology:
compactness, Hausdorffness, and continuity of its group operations
are not required.  The cochains in the conclusion are locally
constant in that topology.
Let \(\kappa:\Delta\to E^\times\) be trivial on \(\pi_U(J)\).
Suppose that the finite local family \(a_v:G_v\to\Delta\)
has the following inertia coverage: for every
\(t\in J\cap I_A\), either \(\pi_U(t)=1\), or
\[
                    \pi_U(t)=g\,a_v(s)g^{-1}
                                                        \tag{LFC.11}
\]
for some \(v,s,g\).  Then the literal strict cone satisfies
\[
 \operatorname{rank}_E H^1_C(a,E(\kappa),0)
 \le \operatorname{clRk}_\ell(J)
           +\#\{v:\kappa a_v=1\}.                       \tag{LFC.12}
\]
This is the original \texttt{strGen}, including its local
primitive correction and arbitrary coefficient field.

\begin{proof}
Restrict a representative \((z,(x_v))\) along \(\pi_U\) to \(J\).
Because \(\kappa\pi_U=1\) there, the restriction is additive.
It is locally constant and is unchanged by adding a
global coboundary.  It kills inertia.  Indeed if (LFC.11)
holds, then \(\kappa(a_v(s))=1\), while strictness says
\(z(a_v(s))=(\kappa(a_v(s))-1)x_v=0\).
The cocycle identity gives the exact formula
\[
 z(gag^{-1})=\kappa(g)z(a)+(1-\kappa(a))z(g),             \tag{LFC.13}
\]
so \(z(\pi_U(t))=0\).  The case \(\pi_U(t)=1\) is immediate.
We have constructed a linear map from the strict cone
space to the \emph{actual} \(X_E(J)\) of LFC7.

To bound its kernel, put \(N=\pi_U(J)\).  This is normal
of finite index \(n\) dividing \([U:J]\), hence
\(\ell\nmid n\).  If the restricted character is zero,
then \(z\) vanishes on \(N\) and descends to
the finite group \(\Delta/N\), as does \(\kappa\).
Set
\[
           X=-n^{-1}\sum_{q\in\Delta/N}z(q).
\]
The cocycle identity and permutation of the summation
indices give \(z(g)=\kappa(g)X-X\).
Subtract the coboundary pair \((dX,(X)_v)\).
The remaining representative has zero global cocycle,
and its local primitive \(x_v-X\) satisfies
\((\kappa(a_v(s))-1)(x_v-X)=0\) for every \(s\).
It is zero whenever \(\kappa a_v\ne1\).
The kernel is therefore an image of
\(E^{\{v:\kappa a_v=1\}}\); it is not silently identified
with zero when a local invariant exists.
Rank--nullity, or the same dimension inequality for image
and kernel, now gives (LFC.12) using LFC7.
\end{proof}
% END LFC8

% BEGIN LFC9
\paragraph{LFC9. The original \texttt{b1p} bound.}
\label{claim:LFC9}
For a full \texttt{FLink}, assume
\[
 \forall v\ \exists s:\ \Psi(\operatorname{loc}_v s)^2\ne1,
 \qquad
 \operatorname{clRk}_\ell
       (K\cap\ker\bar\epsilon_F)\le C_0.                 \tag{LFC.14}
\]
Then
\[
 \operatorname{rank}_k
 H^1_C\bigl(\operatorname{loc},
               k(\Psi^{-2}\delta_k),0\bigr)\le C_0.
                                                        \tag{LFC.15}
\]

\begin{proof}
Apply LFC8 with \(U=G_F\),
\(J=K\cap\ker\bar\epsilon_F\), \(\Delta=\Gamma\), and
\(\kappa=\Psi^{-2}\delta_k\).
The subgroup \(J\) is normal open and its index divides
\([G_F:K](\ell-1)\), so its index is prime to \(\ell\).
Equation (LFC.2) makes \(\kappa\pi\) trivial on \(J\).
For an inertia element at a place outside \(S\), (LFC.1)
kills its image.  At a place in \(S\), conjugacy of
prolongations sends it to the chosen inertia subgroup,
which lies in \(D_{A_v}=\iota_v(G_v)\).  Applying \(\pi\)
gives precisely (LFC.11).
The trivial valuation has trivial inertia and causes
no extra case.

Because \(\delta\operatorname{loc}_v=1\), local triviality
of \(\kappa\) would imply
\(\Psi(\operatorname{loc}_v s)^2=1\) for all \(s\),
contradicting (LFC.14).
The local correction term in (LFC.12) is therefore zero.
Its remaining term is bounded by the same \(C_0\).
\end{proof}
% END LFC9

% BEGIN LFC10
\paragraph{LFC10. The original doubled-local \texttt{b1} bound.}
\label{claim:LFC10}
Under the same hypotheses (LFC.14),
\[
 \operatorname{rank}_k H^1_C
   \bigl(\operatorname{loc}_2,
         k(\Psi^{-1}|_{\Gamma_0}),0\bigr)\le C_0.
                                                        \tag{LFC.16}
\]

\begin{proof}
Use LFC8 with \(U=\pi^{-1}\Gamma_0\),
\(J=K\cap\ker\bar\epsilon_F\), \(\Delta=\Gamma_0\), and
the surjection induced by \(\pi\).  The same index argument
works for \([U:J]\), and \(K\subset U\) by (LFC.2).
The character is trivial on \(J\).
Outside \(S\) the required inertia images vanish.
At \(v\in S\), a conjugating image \(a\in\Gamma\) either
belongs to \(\Gamma_0\), giving the first local branch,
or is \((ac^{-1})c\), giving the second branch with
conjugator \(ac^{-1}\in\Gamma_0\).
Thus the full inertia coverage holds for the actual
family indexed by \(V\times\{0,1\}\).

Characters are conjugation invariant, so
\[
 \Psi(\operatorname{loc}_{2,(v,b)}s)
                  =\Psi(\operatorname{loc}_v s)
                     \quad(b=0,1).
\]
A trivial local \(\Psi^{-1}\) would again contradict the
nonquadratic witness in (LFC.14).
Both branches have zero local correction in (LFC.12),
which proves (LFC.16).  No identification of this doubled
cone with an undoubled cone has been made.
\end{proof}
% END LFC10

\subsection{A shorter proof of the untwisted bound}

% BEGIN LFC11
\paragraph{LFC11. The cyclotomic eigenspace bound over a CM field.}
\label{claim:LFC11}
Let \(F\) be totally real and let \(\ell\) be odd.
Suppose \(\bar\epsilon_F\) is nontrivial on every
decomposition group above \(\ell\).
Put \(E=F(\mu_\ell)\) and let \(r_\ell(E)\) be the
\(\ell\)-rank of its ideal class group.  Then
\[
 \dim_{\mathbb F_\ell}
   \operatorname{Hom}_{\rm cont}
      (G_{F,S},\mathbb F_\ell)\le 1+r_\ell(E),
                                                        \tag{LFC.17}
\]
where \(S\) is the set of \(\ell\)-places.

\begin{proof}
Let \(\Delta=\operatorname{Gal}(E/F)\) and
\(\omega:\Delta\hookrightarrow\mathbb F_\ell^\times\)
be its cyclotomic character.  Its order is prime to \(\ell\).
Restriction of additive characters from \(G_F\) to \(G_E\)
is injective: a character vanishing on \(G_E\) factors
through \(\Delta\), whose order is prime to \(\ell\).
The Kummer construction in LFC1 sends such a restriction
to a class \([a]\in E^\times/E^{\times\ell}\).
For \(d\in\Delta\), its class satisfies
\[
                         d[a]=\omega(d)[a]               \tag{LFC.18}
\]
in additive notation for this \(\mathbb F_\ell\)-space.
Indeed, if \(g\beta/\beta=\zeta^{\chi(g)}\), a lift of \(d\)
sends \(\zeta\) to \(\zeta^{\omega(d)}\), while
\(\chi(d^{-1}gd)=\chi(g)\).  Thus \(d\beta\) is a
Kummer root for the character \(\omega(d)\chi\),
proving (LFC.18) by the injectivity argument of LFC1.
Unramifiedness outside \(\ell\) places the class in
\(V_T(E)/E^{\times\ell}\), with \(T\) the places above \(\ell\).

Taking \(\omega\)-eigenspaces in (LFC.4) is exact:
the projector is
\[
 e_\omega=|\Delta|^{-1}
           \sum_{d\in\Delta}\omega(d)^{-1}d.
 \]
The class-group term has dimension at most \(r_\ell(E)\).
To see this also for the localized class group, it is a
quotient of \(\operatorname{Cl}(\mathcal O_E)\);
the minimal number of generators of a finite abelian
\(\ell\)-group, equivalently the dimension of its
\(\ell\)-torsion, cannot increase on taking a quotient.

We claim that the \(\omega\)-part of the \(T\)-unit term
has dimension at most one.  The valuation map gives
\(\mathcal O_{E,T}^\times/\mathcal O_E^\times\) as a
subgroup \(\Lambda\) of the permutation lattice
\(\mathbb Z[T]\).  The quotient by the units is torsion-free,
so reducing
\(0\to\mathcal O_E^\times\to
 \mathcal O_{E,T}^\times\to\Lambda\to0\)
modulo \(\ell\) is right exact with the left injection.
For a prime \(w\in T\), its stabilizer in \(\Delta\) is
the image of the local decomposition group in
\(\operatorname{Gal}(E/F)\).  By hypothesis \(\omega\)
is nontrivial on this stabilizer.  A vector of the
permutation representation transforming by \(\omega\)
must therefore have zero coefficient at \(w\), and hence
at every prime in its orbit.

Here it is essential to handle the lattice, rather than
assume that reduction preserves an arbitrary inclusion.
Since \(|\Delta|\) is prime to \(\ell\), the idempotent
\(e_\omega\) lifts to \(\mathbb Z_\ell[\Delta]\) using
the Teichmuller lift of \(\omega\).
Its action on \(\mathbb Q_\ell[T]\) is zero by the same
stabilizer argument.  Therefore its action on
\(\Lambda\otimes\mathbb Z_\ell\) is zero, and reduction
shows that \((\Lambda/\ell\Lambda)_\omega=0\).
Thus the \(\omega\)-part of the \(T\)-units modulo
\(\ell\)-th powers comes entirely from ordinary units.

The field \(E=F(\mu_\ell)\) is CM.  Complex conjugation
\(j\in\Delta\) acts by \(\omega(j)=-1\).
For an ordinary unit \(u\), every complex absolute value
of \(u/j(u)\) is one.  This is an algebraic integer unit
all of whose conjugates have absolute value one, hence
a root of unity.  One elementary proof of this last fact
observes that the monic polynomials having the conjugates
of its positive powers as roots have integral coefficients
bounded uniformly by binomial coefficients.  There are
only finitely many such polynomials and thus finitely many
possible powers; two powers agree.
Consequently \((1-j)\mathcal O_E^\times\subset\mu(E)\).
Modulo \(\ell\), every \(j\)-anti-invariant unit class is
half of its image under \(1-j\), and so lies in the image
of \(\mu(E)/\mu(E)^\ell\).
The roots of unity in a number field form a finite cyclic
group, whose quotient by \(\ell\)-th powers has dimension
at most one.  This proves the claim about units.
Combining it with the class-group bound proves (LFC.17).
\end{proof}
% END LFC11

% BEGIN LFC12
\paragraph{LFC12. The full original \texttt{h1} field.}
\label{claim:LFC12}
For a full \texttt{FLink}, suppose \(k\) is finite,
\(\omega_k\operatorname{loc}_v\ne1\) for every \(v\), and
\(\operatorname{clRk}_\ell(K\cap\ker\bar\epsilon_F)\le C_0\).
Then
\[
 \operatorname{rank}_k
       H^1_C(\operatorname{loc},k(1),k)\le1+C_0.          \tag{LFC.19}
\]
In particular the source's theorem \texttt{GC.h1}, with
its additional retained premise \(P\), follows.

\begin{proof}
With the full local subspaces \(k\), the local primitives
in (LFC.3) are quotiented out.  Since the action is trivial,
global coboundaries are zero.  Thus the space in (LFC.19)
is literally \(\operatorname{Hom}_{\rm cont}(\Gamma,k)\).
For an additive \(\mathbb F_\ell\)-character of \(\Gamma\),
its pullback to \(G_F\) kills every inertia outside \(S\):
for \(t\in I_A\), the prime-to-\(\ell\) power
\(t^{[G_F:K]}\) belongs to \(I_A\cap K\), and (LFC.1)
kills its image.  Hence LFC11 bounds the prime-field
character dimension by \(1+r_\ell(F(\mu_\ell))\).

Put \(J=K\cap\ker\bar\epsilon_F\) and \(E_1=\Omega^J\).
Then \(E_1\supset E=F(\mu_\ell)\) and \([E_1:E]\)
is prime to \(\ell\).
Restriction of unramified additive characters
\(G_E\to\mathbb F_\ell\) to \(G_{E_1}\) is injective.
Indeed, for a character vanishing on the latter subgroup,
the \([E_1:E]\)-th power of every element lies in it
(these are normal subgroups here), and that degree is
invertible modulo \(\ell\).
Restriction remains unramified because inertia is
intersected with \(G_{E_1}\).
Finite class-field reciprocity and LFC7 therefore give
\[
 r_\ell(E)\le\dim_{\mathbb F_\ell}X_{\mathbb F_\ell}(J)
             =\operatorname{clRk}_\ell(J)\le C_0.
 \]
Finally a basis of the finite extension \(k/\mathbb F_\ell\)
shows, by taking coordinates of additive functions, that
\(\operatorname{Hom}_{\rm cont}(\Gamma,k)\) is obtained
from its prime-field character space by scalar extension.
Its dimension over \(k\) is the prime-field dimension.
This gives (LFC.19) with the original \(C_0\).
\end{proof}
% END LFC12

\subsection{Euler estimates through the actual inertia quotient}

% BEGIN LFC13
\paragraph{LFC13. A finite ramification level with its exact kernel.}
\label{claim:LFC13}
For a \texttt{GLink} over any number field there is a
finite set \(S'\supset S\)
and a continuous surjection
\[
 p:G_{F,S'}\twoheadrightarrow\Gamma,\qquad p\,q=\pi,
                                                        \tag{LFC.20}
\]
where \(q:G_F\to G_{F,S'}\), such that
\[
 \ker p=\overline{\langle J_v:v\in S'\setminus S
                                      \rangle^{\rm normal}},
 \qquad J_v=q(I_{A_v}\cap K).                            \tag{LFC.21}
\]
One chosen prolongation at each place suffices in this formula.

\begin{proof}
The finite extension fixed by \(K\) ramifies at finitely
many finite places, say \(T\); this follows by inverting
the nonzero discriminant of an integral basis.
At \(v\notin T\) all inertia is contained in \(K\).
Take \(S'=S\cup T\).
For \(v\notin S'\), (LFC.1) therefore kills all of \(I_A\),
so \(\pi\) factors through the maximal quotient
unramified outside \(S'\), giving (LFC.20).
The quotient topology gives continuity, and surjectivity
comes from that of \(\pi\).

Under \(q\), all generators of (LFC.1) at places outside
\(S'\) become one.  At a place of \(S'\setminus S\),
any prolongation is conjugate to the chosen one, and
normality of \(K\) carries its inertia intersection to
the chosen intersection.  Thus their images are precisely
the normal generators in (LFC.21).
To justify closure, the image of a closed subgroup of a
compact group in a Hausdorff quotient is closed.
The closed normal subgroup generated by the displayed
images has inverse image containing every generator
of \(\ker\pi\); conversely its generators all lie in
\(\ker p\).  These two inclusions prove (LFC.21).
\end{proof}
% END LFC13

% BEGIN LFC14
\paragraph{LFC14. The exact local correction at a discarded prime.}
\label{claim:LFC14}
Here \(\ell\) may be any prime, including two, and the
base number field is arbitrary.
Let \(v\nmid\ell\), let \(D=D_A\), \(I=I_A\), and
\(J=I\cap K\).  Let \(W\) be a finite-dimensional
representation over a finite field \(E\) of characteristic
\(\ell\), with continuous action and trivial \(J\)-action.
Define
\[
 \mathcal E(D,J;W)=
 \{f:J\to W:\ f\text{ continuous and additive},\
                     f(g t g^{-1})=\rho(g)f(t)\}.
 \]
Under the local fields of \(P\), this is finite-dimensional and
\[
              \dim_E\mathcal E(D,J;W)=\dim_E H^2(D,W).
                                                        \tag{LFC.22}
\]
Moreover every class inflated to \(D\) from a continuous
quotient killing \(J\) is zero in \(H^2(D,W)\).

\begin{proof}
Here \([I:J]\) divides \([G_F:K]\), so LFC4 applies.
Restriction of a one-cocycle to \(J\) is an additive
equivariant map and kills coboundaries.
Its kernel in \(H^1(D,W)\) is \(H^1(D/J,W)\):
if the restriction is zero, the cocycle descends
continuously through the compact-to-Hausdorff quotient;
the converse is immediate.

We verify surjectivity of this restriction without
assuming an extension property for arbitrary functions.
For \(f\in\mathcal E(D,J;W)\), choose a normal open
subgroup \(U_0\subset D\) killing the action and satisfying
\(U_0\cap J\subset\ker f\).  Such a subgroup exists
by continuity and compactness.  On the finite quotient
by \(U_0\), choose representatives modulo the image of \(J\)
and define a function \(a:D\to W\) by
\[
                 a(r t)=\rho(r)f(t),\qquad t\in J.
 \]
The condition on \(U_0\cap J\) makes this definition
well defined on the quotient; choose value zero on
the representatives.  It is continuous and satisfies
\(a(g t)=a(g)+\rho(g)f(t)\).
Equivariance of \(f\) shows by substitution that
\(d a(g,h)\) is unchanged on replacing either \(g\) or \(h\)
by a right translate by \(J\).  It is therefore a
two-cocycle on \(D/J\).
By LFC4 it is \(d b\) for a continuous cochain \(b\)
on \(D/J\).  The cocycle \(a-b\) restricts to \(f\).
This proves the exact sequence
\[
 0\longrightarrow H^1(D/J,W)\longrightarrow H^1(D,W)
       \longrightarrow\mathcal E(D,J;W)\longrightarrow0.
 \]
LFC4 gives dimension \(\dim H^0(D,W)\) for its first term.
The local Euler characteristic away from \(\ell\) is
\(\dim H^1(D,W)=\dim H^0(D,W)+\dim H^2(D,W)\).
This proves (LFC.22), with finiteness supplied by \(P\).
For the final assertion, the quotient map from \(D\)
factors through \(D/J\), whose \(H^2\) is zero by LFC4.
The pullback of every two-cocycle class is therefore zero.
\end{proof}
% END LFC14

% BEGIN LFC15
\paragraph{LFC15. The inflation defect is bounded by the same corrections.}
\label{claim:LFC15}
Let \(p:\Gamma'\twoheadrightarrow\Gamma\) be a continuous
surjection of profinite groups whose kernel is the closed
normal subgroup generated by a finite family \(J_v\).
Let \(D_v\) be arbitrary subgroups of \(\Gamma'\), and let
\(W\) be a finite representation of \(\Gamma\) over a finite
field of any prime characteristic.
No containment \(J_v\subset D_v\) or normalization of \(J_v\)
by \(D_v\) is assumed.  Here \(\mathcal E(D_v,J_v;W)\)
means the continuous additive functions on \(J_v\)
satisfying
\(f(g t g^{-1})=\rho(p g)f(t)\) whenever \(g\in D_v\),
\(t\in J_v\), and \(g t g^{-1}\in J_v\).
This is the source's conditional equivariance convention.
Put
\[
 T=\ker\bigl(H^2(\Gamma,W)\to H^2(\Gamma',W)\bigr).
 \]
If \(H^1(\Gamma',W)\) and the spaces
\(\mathcal E(D_v,J_v;W)\) are finite-dimensional, then
\(H^1(\Gamma,W)\) and \(T\) are finite-dimensional and
\[
 \dim T+\dim H^1(\Gamma',W)
   \le\dim H^1(\Gamma,W)+
                          \sum_v\dim\mathcal E(D_v,J_v;W).
                                                        \tag{LFC.23}
\]

\begin{proof}
Inflation on \(H^1\) is injective: a pulled-back
cocycle which is \(g\mapsto\rho(p g)X-X\) has the
same formula downstairs, by surjectivity of \(p\).
Restriction to the \(J_v\)'s defines
\[
 R:H^1(\Gamma',W)\longrightarrow
                  \bigoplus_v\mathcal E(D_v,J_v;W).
 \]
Its kernel is the image of inflation.
Indeed a cocycle vanishing on every \(J_v\) vanishes
on their conjugates by the cocycle formula, and its zero
set inside \(\ker p\) is a closed normal subgroup.
It thus vanishes on \(\ker p\) and descends continuously.

For the second-degree defect, consider the vector space
\(\mathcal T\) of pairs \((z,\beta)\), where \(z\) is a
normalized continuous two-cocycle on \(\Gamma\),
\(\beta(1)=0\), and \(d\beta=p^*z\).
It maps onto \(T\) by \((z,\beta)\mapsto[z]\).
Every element of \(T\) has such a representative:
subtract the coboundary of the constant value \(\beta(1)\)
to normalize a primitive and the corresponding cocycle.
For \(n\in\ker p\), the identity \(d\beta=p^*z\) gives
\[
 \beta(ng)=\beta(n)+\beta(g),\quad
 \beta(gn)=\beta(g)+\rho(p g)\beta(n),\quad
 \beta(gng^{-1})=\rho(p g)\beta(n).                       \tag{LFC.24}
\]
Hence restrictions of \(\beta\) define a linear map
\(r:\mathcal T\to\bigoplus_v\mathcal E(D_v,J_v;W)\).
If \(r(z,\beta)=0\), (LFC.24) and closed normal generation
force \(\beta|_{\ker p}=0\).
The first two identities then make \(\beta\) constant
on fibers, so it descends to a continuous \(b:\Gamma\to W\).
Surjectivity gives \(z=db\), and \([z]=0\).
Thus \(\ker r\) is contained in the kernel of
\(\mathcal T\to T\).
The restriction image \(R(H^1(\Gamma',W))\) is contained
in \(r(\ker(\mathcal T\to T))\), using the pairs \((0,\beta)\).
It follows that \(T\) is a quotient of a subspace of
\[
 \left(\bigoplus_v\mathcal E(D_v,J_v;W)\right)
                    \big/R(H^1(\Gamma',W)).
 \]
In particular it is finite-dimensional.
Since
\(\dim R(H^1(\Gamma',W))=
 \dim H^1(\Gamma',W)-\dim H^1(\Gamma,W)\),
the displayed quotient bound is exactly (LFC.23).
This argument uses actual continuous primitives and
requires no global continuous section of \(p\).
\end{proof}
% END LFC15

% BEGIN LFC16
\paragraph{LFC16. The finite-level Euler and Poitou--Tate identities.}
\label{claim:LFC16}
Let \(F\) be totally real, \(\ell\) odd,
\(S'\) finite containing every \(\ell\)-place, and \(W\)
a finite-dimensional continuous representation of
\(G_{F,S'}\) over a finite field \(E\) of characteristic
\(\ell\).  The retained fields of \(P\) give finite
\(H^1,H^2\) and the identities
\[
 \begin{split}
 \dim H^1+\sum_{w\mid\infty}\dim W^{c_w=1}
   &=\dim H^0+\dim H^2+[F:\mathbb Q]\dim W,\\
 \dim\operatorname{Sha}^2+\sum_{v\in S'}\dim H^2(D_v,W)
   &=\dim H^2+\dim H^0(G_F,W^\vee(1)).
 \end{split}                                           \tag{LFC.25}
\]
Here \(\operatorname{Sha}^2\) is precisely the kernel of localization
to the displayed \(H^2(D_v,W)\)'s.

\begin{proof}
Apply \(P\)'s local and global finiteness and its global
Euler characteristic at exponent \(\ell^1\), after
restricting scalars to \(\mathbb F_\ell\).
The cocycle and coboundary sets do not change under this
restriction: their defining equations and continuity
conditions are identical.  If \(q=|E|\), each finite
cohomology space therefore has cardinality \(q^{\dim_E}\).
The Euler characteristic's cardinal identity becomes
the first equality after taking exponents of \(q>1\).

The last two maps in the normalized Poitou--Tate sequence
provided by \(P\) are exact, and the last is onto:
\[
 H^2(G_{F,S'},W)\longrightarrow
 \bigoplus_{v\in S'}H^2(D_v,W)
       \longrightarrow H^0(G_F,W^\vee(1))^\vee
       \longrightarrow0.
 \]
Its first kernel is the stated \(\operatorname{Sha}^2\), by the
literal restriction maps on two-cocycles.
Taking cardinalities gives the second identity.
The coefficient-field dual is correctly identified as
follows.  Composition with
\(\operatorname{Tr}_{E/\mathbb F_\ell}\) identifies the
underlying additive space of \(\operatorname{Hom}_E(W,E)\)
with \(\operatorname{Hom}_{\mathbb F_\ell}(W,\mathbb F_\ell)\).
Injectivity follows from nondegeneracy of the finite-field
trace pairing, and equal cardinalities give surjectivity.
The map commutes with the representation and the
cyclotomic twist, so it identifies the invariant spaces.
This supplies the exact dual term in (LFC.25), without
assuming an unintended \(E\)-linearity of the original
prime-field duality normalization.
\end{proof}
% END LFC16

% BEGIN LFC17
\paragraph{LFC17. The global estimate for the actual quotient.}
\label{claim:LFC17}
For a \texttt{GLink}, a totally real \(F\), and a
finite-dimensional finite-field representation \(W\)
of \(\Gamma\) with continuous action, \(H^1(\Gamma,W)\)
and \(H^2(\Gamma,W)\) are finite-dimensional, and
\[
 \begin{split}
 &\dim\operatorname{Sha}^2_S(\Gamma,W)+
       \sum_{v\in S}\dim H^2(D_v,W)+\dim H^0(\Gamma,W)
                          +[F:\mathbb Q]\dim W\\
 &\hspace{12mm}\le
       \dim H^1(\Gamma,W)+
       \sum_{w\mid\infty}\dim W^{\pi(c_w)=1}
       +\dim H^0(G_F,W^\vee(1)).
 \end{split}                                           \tag{LFC.26}
\]
The local actions on the left are through \(\pi\), and
\(\operatorname{Sha}^2_S\) is their exact localization kernel.

\begin{proof}
Choose \(S',p\) by LFC13.
Let \(a_i=\dim H^i(\Gamma,W)\),
\(a_i'=\dim H^i(G_{F,S'},W)\),
\(t=\dim\ker(H^2(\Gamma,W)\to H^2(G_{F,S'},W))\),
and \(b_v=\dim H^2(D_v,W)\).
At a discarded place \(v\in S'\setminus S\), pullback
injects the equivariant homomorphisms on
\((qD_v,q(I_v\cap K))\) into those on \((D_v,I_v\cap K)\).
LFC14 bounds their dimension by \(b_v\).
LFC15 consequently gives
\[
                     t+a_1'\le a_1+\sum_{S'\setminus S}b_v.
                                                        \tag{LFC.27}
\]
It also gives finiteness of \(a_1\) and \(t\).
The target \(H^2(G_{F,S'},W)\) is finite by LFC16;
its finite kernel \(T\) then proves finiteness of \(a_2\).

Inflation sends \(\operatorname{Sha}^2_S(\Gamma,W)\) into
\(\operatorname{Sha}^2_{S'}(G_{F,S'},W)\).
At places in \(S\) this is the definition; at the other
places LFC14 says that all inflated classes vanish
locally, because \(\pi\) kills \(I_v\cap K\).
Thus
\[
             \dim\operatorname{Sha}^2_S\le t+\dim\operatorname{Sha}^2_{S'}.
 \]
Use the second equality in (LFC.25) to replace
\(\dim\operatorname{Sha}^2_{S'}\), and then the first to replace
\(a_2'\).  The invariants and the real fixed spaces
are unchanged by the surjection \(p\).
The resulting discarded-place sum is cancelled by
(LFC.27).  Explicitly the left side of (LFC.26) is at most
\[
 t+a_2'+\dim H^0(G_F,W^\vee(1))
     -\sum_{S'\setminus S}b_v+a_0+[F:\mathbb Q]\dim W,
 \]
which is at most its right side by these two identities.
This proves the estimate without erasing any discarded
prime or assuming inflation on \(H^2\) is injective.
\end{proof}
% END LFC17

% BEGIN LFC18
\paragraph{LFC18. Local coefficient quotients preserve the required bound.}
\label{claim:LFC18}
In LFC17 suppose \(H^0(G_F,W^\vee(1))=0\).
For every family of equivariant surjections
\(q_v:W\twoheadrightarrow Q_v\) for the local groups,
let
\[
 X_Q=\ker\left(H^2(\Gamma,W)\longrightarrow
                              \bigoplus_{v\in S}H^2(D_v,Q_v)\right).
 \]
Then
\[
 \begin{split}
 \dim X_Q+\sum_{v\in S}\dim H^2(D_v,Q_v)
       +\dim H^0(\Gamma,W)+[F:\mathbb Q]\dim W\\
 \le \dim H^1(\Gamma,W)
       +\sum_{w\mid\infty}\dim W^{\pi(c_w)=1}.
 \end{split}                                           \tag{LFC.28}
\]
No flatness, choice of a local complement preserved by
the action, or global stability of \(\ker q_v\) is required.

\begin{proof}
Local Tate duality identifies \(H^2(D_v,W)\) with the
dual of \(\operatorname{Hom}_{D_v}(W,E(1))\), and similarly
for \(Q_v\).  These perfect finite pairings are natural
in the coefficients: on cochains this is the equality
\(f(q_v x)=(f\circ q_v)(x)\) in the evaluation pairing.
Thus the map on \(H^2\) induced by \(q_v\) is dual to the
injection
\(\operatorname{Hom}_{D_v}(Q_v,E(1))
 \hookrightarrow\operatorname{Hom}_{D_v}(W,E(1))\).
It is therefore onto.  The same conclusion follows after
the trace identification in LFC16, so it holds with the
precise prime-field duality contained in \(P\).

Here is the needed linear algebra.  For maps
\(r_v:A\to B_v\) and surjections \(f_v:B_v\to C_v\),
put \(X=\bigcap\ker(f_vr_v)\) and \(Y=\bigcap\ker r_v\).
Restriction of \((r_v)_v\) gives an injection
\(X/Y\hookrightarrow\bigoplus_v\ker f_v\).
Hence
\[
          \dim X+\sum_v\dim C_v
                   \le\dim Y+\sum_v\dim B_v.
 \]
Apply this to \(A=H^2(\Gamma,W)\), the local
\(H^2(D_v,W)\), and their quotient-coefficient images.
LFC17 with its zero dual term gives (LFC.28).
\end{proof}
% END LFC18

\subsection{From ordinary cohomology to the literal cone}

% BEGIN LFC19
\paragraph{LFC19. The full cone dimension identity.}
\label{claim:LFC19}
This claim is independent of the arithmetic linked data.
Let \(\Gamma\) and \(G_v\) be arbitrary profinite groups,
with finitely many \(v\), continuous maps
\(\operatorname{loc}_v:G_v\to\Gamma\), a finite field \(E\),
and a finite-dimensional continuous representation \(W\) of
\(\Gamma\) over \(E\).
Let \(q_v:W\twoheadrightarrow Q_v=W/b_v\) be equivariant
for each local group, and suppose the global \(H^1,H^2\)
and local \(H^1(G_v,Q_v)\) are finite-dimensional.
Equivalently one may start with arbitrary equivariant
surjections \(q_v\), taking \(b_v=\ker q_v\) and identifying
the targets with \(W/b_v\); all the maps below respect
that canonical identification.
Write
\[
 \begin{split}
 H^0_C&=\{x\in W^\Gamma:x\in b_v\text{ for every }v\},\\
 X_Q&=\bigcap_v\ker\bigl(H^2(\Gamma,W)\to H^2(G_v,Q_v)\bigr).
 \end{split}
 \]
Then the literal \(H^1_C,H^2_C\) are finite-dimensional and
\[
 \begin{split}
 &\dim H^2_C+\dim H^0_C+\dim H^1(\Gamma,W)
                       +\sum_v\dim H^0(G_v,Q_v)\\
 &\qquad=\dim H^1_C+\dim H^0(\Gamma,W)
                       +\sum_v\dim H^1(G_v,Q_v)+\dim X_Q.
 \end{split}                                           \tag{LFC.29}
\]

\begin{proof}
We use the cochain model, so that the local primitives
are accounted for rather than discarded.
A degree-two cone representative is a global continuous
two-cocycle \(z\), together with local continuous
one-cochains \(a_v:G_v\to W\), satisfying
\[
 q_v z(\operatorname{loc}_v s,\operatorname{loc}_v t)
                         =d(q_v a_v)(s,t).
 \]
Its coboundaries have the form
\[
 z=d\beta,\qquad
 a_v(s)\equiv
   \beta(\operatorname{loc}_v s)
        -\bigl(\rho(\operatorname{loc}_v s)Y_v-Y_v\bigr)
                                    \pmod {b_v}.        \tag{LFC.30}
\]
These are exactly the source's
\texttt{continuousCocycles2} and
\texttt{continuousCoboundaries2}.
Choose arbitrary linear sections of the \(q_v\)'s;
they are continuous for the discrete topologies and
are not required to intertwine the action.

The definitions give the following exact sequence, ending
with the indicated image:
\[
 \begin{gathered}
 0\longrightarrow H^0_C\longrightarrow H^0(\Gamma,W)
 \longrightarrow \bigoplus_v H^0(G_v,Q_v)
 \longrightarrow H^1_C\longrightarrow H^1(\Gamma,W)\\
 \longrightarrow\bigoplus_v H^1(G_v,Q_v)
 \longrightarrow H^2_C\longrightarrow X_Q
 \longrightarrow0.                                    \tag{LFC.31}
 \end{gathered}
\]
Here an invariant local tuple maps to the class of
\((0,(x_v))\) after lifting it; a local one-cocycle
tuple maps to the cone class \((0,(a_v))\) after lifting.
The maps from the cone groups forget the local data.
Changing a lift changes its value by \(b_v\), hence
does not change these classes.

For clarity, exactness can be checked directly at every
nontrivial step.  A local invariant tuple is a cone
coboundary precisely when it is the image of one global
invariant \(X\), by (LFC.3).
If the global cocycle of a cone-one class is \(dX\),
subtract its coboundary; the remaining local primitives
are invariant in \(Q_v\), giving the preceding image.
A global one-class has zero local classes precisely
when its local coboundary primitives can be lifted to
form a cone-one representative.
A local one-cocycle tuple gives zero in \(H^2_C\)
precisely when (LFC.30) supplies a global one-cocycle
\(\beta\) whose local cohomology classes equal that tuple.
If a cone-two class has zero global class, subtract
the coboundary of a global primitive \(\beta\);
the remaining \(q_v(a_v-\beta\operatorname{loc}_v)\)
are local one-cocycles.
Finally a global two-class lies in \(X_Q\) precisely
when local continuous primitives exist in each \(Q_v\);
lift them to obtain a cone-two representative.
This proves (LFC.31).

Every term except possibly the cone terms is
finite-dimensional; exactness proves their finiteness
as well.  The alternating dimension sum of (LFC.31)
is zero, and rearranging it is exactly (LFC.29).
\end{proof}
% END LFC19

% BEGIN LFC20
\paragraph{LFC20. The complete \texttt{EC} field for every finite module.}
\label{claim:LFC20}
Assume \(P\) and a full \texttt{FLink}.
For every finite field \(E\) of characteristic \(\ell\),
finite-dimensional \(E\)-space \(W\), representation
\(\rho:\Gamma\to\operatorname{End}_E(W)\) with open kernel,
and locally stable subspaces \(b_v\subset W\), suppose
\[
 \bigl[f(\rho(g)x)=\omega_E(g)f(x)
           \text{ for all }g,x\bigr]\Longrightarrow f=0
                   \quad(f:W\to E\text{ linear}).        \tag{LFC.32}
\]
Then \(H^2_C\) is finite-dimensional and
\[
 \dim_E H^2_C+\dim_E H^0_C+\sum_v d_v\dim_E b_v
 \le \dim_E H^1_C+
                 \sum_{w\in\mathrm{Inf}}\dim_E W^{c_w=1}.
                                                        \tag{LFC.33}
\]
This is the full source \(\Phi.\texttt{EC}\), with its
quantifiers over coefficient fields, representations,
and all stable local conditions.

\begin{proof}
An open kernel makes the action continuous: for fixed
\(x\), its value is constant on cosets of that kernel;
the finite discrete space \(W\) then gives joint
continuity.  Condition (LFC.32), surjectivity of \(\pi\),
and \(\epsilon\pi=\epsilon_F\) say exactly that
\(H^0(G_F,W^\vee(1))=0\).
Put \(Q_v=W/b_v\).
The maps \(\iota_v:G_v\simeq D_{A_v}\) are homeomorphisms,
since they are continuous bijections from compact
groups to Hausdorff groups.  Pullback and its inverse
identify their continuous cocycles, coboundaries,
invariants and cohomology.  Thus LFC18 applies to
the actual local family of the frame.

Local Euler characteristic at \(v\mid\ell\) gives
\[
 \dim H^1(G_v,Q_v)+d_v\dim b_v
 =\dim H^0(G_v,Q_v)+\dim H^2(G_v,Q_v)+d_v\dim W.
                                                        \tag{LFC.34}
\]
All local groups in this equation are finite-dimensional.
The local degree formula, with every place above \(\ell\)
present exactly once, gives
\(\sum_v d_v=[F:\mathbb Q]\).
Use (LFC.34) in (LFC.29).
The combination on its right involving \(X_Q\), local
\(H^2\), global \(H^0\), and \([F:\mathbb Q]\dim W\)
is bounded by (LFC.28).
Cancel the remaining global \(H^1\) and local \(H^0\)
terms.  This is precisely (LFC.33), with the real
fixed spaces identified by the given bijection
\(\mathrm{Inf}\simeq\{\text{real places of }F\}\).
Finiteness of \(H^2_C\) follows from LFC17 and LFC19.
In particular no chosen local condition was replaced
by \(0\) or by \(W\) in proving this field.
\end{proof}
% END LFC20

\subsection{Assembly without changing the field, marking, or bound}

% BEGIN LFC21
\paragraph{LFC21. All twenty-eight \texttt{Ok} fields on the MFC frame.}
\label{claim:LFC21}
Fix the full original \(\texttt{Gen}\), \(P\), and the
single \(L,g_1,\Phi,C_0,w\) constructed in MFC14.
Then this very frame satisfies
\[
                         \Phi.\texttt{Ok}(C_0)
                    \quad\hbox{and}\quad
                         \Phi.\texttt{Big}(C_0).          \tag{LFC.35}
\]

\begin{proof}
Let \(F=L^+\), \(U=G_F\subset G_{\mathbb Q}\), and
\(\alpha=\overline\Psi_{\mathbb Q}\).
The \texttt{FLink} supplied by the \texttt{Linked} witness
has
\[
 K=(\ker\eta\cap\ker\alpha)|_{G_F},\qquad
 J=K\cap\ker\bar\epsilon_F.
 \]
Under the actual field-embedding map, \(J\) is precisely
\[
     U\cap\ker\eta\cap\ker\alpha\cap\ker\bar\epsilon_{\mathbb Q}.
                                                        \tag{LFC.36}
\]
This is an equality of subgroups after the natural
topological isomorphism, not just an index comparison.
Inertia groups are carried to intersections with the
same valuation inertia, and the cyclotomic characters
agree on every root of unity.  Precomposition therefore
identifies the literal \(\texttt{XDual}(1,-,\mathbb F_\ell,
\mathbb F_\ell)\) spaces, including continuity and
every inertia condition.
Their ranks are equal.  Thus MFC5's uniform bound,
at MFC's already chosen layer \(m\), is exactly
\(\operatorname{clRk}_\ell(J)\le C_0\) required above.

MFC4 and MFC8 give a nontrivial local
\(\bar\epsilon_F\), at every local group, and MFC9
gives the nonquadratic residual local witness
\(\Psi(\operatorname{loc}_v s)^2\ne1\).
The conversion of cyclotomic reduction to
\(\omega_k\) uses the injective prime-field map,
as in LFC6, and does not strengthen a residue condition.
Apply LFC2 for \texttt{phi}, LFC6 for \texttt{gen},
LFC3 for \texttt{cl}, LFC20 for \texttt{ec},
LFC9 for \texttt{b1p}, LFC10 for \texttt{b1},
and LFC12 for \texttt{h1}, using this one \(C_0\).

The twenty-one already proved fields are
\[
 \begin{gathered}
 \texttt{l\_odd},\ \texttt{fin},\ \texttt{del\_open},\
 \texttt{c\_sq},\ \texttt{del\_c},\ \texttt{loc\_cont},\
 \texttt{eps\_cont},\ \texttt{Psi\_open},\ \texttt{Psi\_c},\\
 \texttt{baseSubgroup\_spec},\ \texttt{card\_inf},\
 \texttt{cw\_sq},\ \texttt{del\_cw},\
 \texttt{f2a},\ \texttt{f2b},\ \texttt{f2c},\
 \texttt{f3},\ \texttt{f4},\ \texttt{f4'},\
 \texttt{f5},\ \texttt{Psi\_N}.
 \end{gathered}
 \]
Here the displayed \texttt{Psi} is the source's
\(\Psi\)-named field.  MFC13 supplies both nonzero
strict cone classes for \texttt{f4}, \texttt{f4'} on
this same quotient.  Together with the seven just
proved, these are exactly all twenty-eight fields.
MFC9 and MFC14 already proved
\(3C_0+7<\sum_v d_v\) and \(d_v>5\) for every \(v\),
which are the two \texttt{Big} requirements.
They are unchanged.  This proves (LFC.35).
\end{proof}
% END LFC21

% BEGIN LFC22
\paragraph{LFC22. The full original \texttt{HasFrameW} constructor.}
\label{claim:LFC22}
For every commutative local ring \(O\) whose residue
field has characteristic \(\ell\), the full original
\(\texttt{Gen}(\ell,\eta,\Psi,\sigma,B,M,H_b)\) and \(P\)
imply, for every \(Q>0\),
\[
                \texttt{HasFrameW}
                     (\ell,\eta,\Psi,\sigma,B,M,Q).
                                                        \tag{LFC.37}
\]
The output is one CM field \(L\subset\Omega\), one
marked \(g_1\in G_L\cap D_B\), one frame \(\Phi\),
and \(C_0,w\), with all the literal output clauses.

\begin{proof}
Fix \(Q\) first.  MFC1 chooses the real quadratic
subgroup \(H_0\).  CAF5 chooses a witness prime \(w\)
with \(w\equiv1\pmod Q\); its witness conditions are
the complete \texttt{FrobeniusWitnessPrime}.
MFC5 chooses \(C_0\) uniformly for every later layer,
using TW1 at the exact fixed abelian field.
Only then set \(m=3C_0+7\), so
\(3C_0+7<\ell^m\).
MFC4, MFC6 and CAF8 construct the actual CM fixed
field with
\[
  U=H_0\cap H_\theta(w)\cap H_b\cap H_m,\qquad
  G_L=U\cap\ker\eta,\qquad G_{L^+}=U.
 \]
The nonquadratic local element provided there is the
marked \(g_1\in D_B\cap G_L\).
MFC8 constructs \(\Phi\) and the full \texttt{Linked}
witness for this \(L,g_1\); MFC12--MFC13 supply both
strict classes, and LFC21 gives its full \texttt{Ok}
and \texttt{Big} at the already chosen \(C_0\).
Finally \(G_{L^+}=U\subset H_\theta(w)\) follows directly
from the displayed intersection.

Thus the same tuple has, in the order of the definition:
\(\texttt{Linked}\), \(\Phi.\texttt{Ok}(C_0)\),
\(\Phi.\texttt{Big}(C_0)\), the full witness-prime record,
\(w\equiv1\pmod Q\), and
\(\operatorname{upl}(\Omega,L)\subset H_\theta(w)\).
These are exactly all the fields of \texttt{HasFrameW}.
The finite residue-field requirement was supplied by
\texttt{Gen.fin}; no DVR or characteristic-zero condition
on \(O\) was used for (LFC.37).
\end{proof}
% END LFC22

% BEGIN LFC23
\paragraph{LFC23. The original coefficient-DVR \texttt{HasFrame} join.}
\label{claim:LFC23}
Specialize to the original coefficient DVR \(O=c.O\)
and retain the full \texttt{Gen} and \(P\).
The construction gives \(\texttt{HasFrame}\), with
the character \(\texttt{SeedClause}\) on the same
\(L,\Phi,g_1\) as its complete \texttt{Ok} and
\texttt{Big} fields.

\begin{proof}
CAF7, applied to the full \texttt{Gen}'s seed-relevant
fields in this coefficient-DVR setting, chooses one
positive integer \(Q\) and the core finite characters
before the later compatible \(L\) and \(w\).
Apply LFC22 to this \(Q\).  The produced \(w\) has
the stipulated congruence, and the resulting \(L\)
has the required embedded fixed-field description
and \(\operatorname{upl}(\Omega,L)\subset H_\theta(w)\).
CAF7 therefore supplies the complete original
\texttt{SeedClause} for that very field.
CAF11's literal installation then returns the
\texttt{GalLink}, the complete \(\texttt{Ok}\),
\(\texttt{Big}\), and this \texttt{SeedClause} with
the same marked element and frame.
This is the defining conjunction of \texttt{HasFrame}.
The statement is restricted to \(O=c.O\) because that
is CAF7's original coefficient-DVR scope.
Its character seed is not an automorphic rank-two
or rank-three representation.
\end{proof}
% END LFC23

\paragraph{Exact remaining arithmetic boundary.}
The seven missing \texttt{Ok} fields, the full
\texttt{frame\_ex}/\texttt{HasFrameW} theorem, and
the coefficient-DVR \texttt{HasFrame} join have now
been proved at their original premises.
To obtain \texttt{GdFrameEx} and \texttt{HostFrameEx}
from their named data, one must still supply the full
entry records \texttt{gen\_gd} and \texttt{gen\_host}.
The latter requires the actual
\(H_b=\ker(\operatorname{dirG}(\kappa_B))\), its
\(\ell\)-power index, its inertia containment and exact
Teichmuller lifting, as well as the host's local
nonquadratic witness.  The host case continues to
allow \(\ell=3\); no stronger lower bound on \(\ell\)
has been introduced in this chapter.
Both automorphic seed producers, the shared coefficient
extension, the original/alternative endpoint matching,
and the full common \texttt{GdEnd} remain separate work.
No whole selected compound obligation or complete
human proof of RH or Dirichlet GRH is closed by this
chapter alone.
